% Generated by scripts/assemble.py; corrections belong in source/corrections.json.
\examhead{128}{郑州大学 2025 年数学分析试卷}{122}
% p123-b01; source PDF page 123
\blockanchor{p123-b01}
\question{1}(10 分) 求极限 $\lim_{x\to 0}\left[\frac{1 + \int_0^x\sin t\mathrm{d}t}{\ln(1 + x^2)} -\frac{1}{\sin^2x}\right].$\par
% p123-b02; source PDF page 123
\blockanchor{p123-b02}
\question{2}(10 分) 求极限 $\lim_{n\to \infty}\frac{(1 + 2^a + 3^a + \cdots + n^a)^{b + 1}}{[1 + 3^b + 5^b + \cdots + (2n + 1)^b]^{a + 1}}$，其中 $a > -1, b > -1$.\par
% p123-b03; source PDF page 123
\blockanchor{p123-b03}
\question{3}(10 分) 求不定积分 $\int \frac{x^2}{(1 - x)^{2024}} \mathrm{d}x$.\par
% p123-b04; source PDF page 123
\blockanchor{p123-b04}
\question{4}(10 分) 判断级数 $\sum_{n=1}^{\infty} (-1)^{n-1} \left[ e - \left( 1 + \frac{1}{n} \right)^n \right]$ 的绝对收敛性和条件收敛性.\par
% p123-b05; source PDF page 123
\blockanchor{p123-b05}
\question{5}(10 分) 设 $f(x) = \arctan 2x$, 求 $f^{(n)}(0)$.\par
% p123-b06; source PDF page 123
\blockanchor{p123-b06}
\question{6}(15 分) 设 $D = [0, 1] \times [0, 1]$.\par
% p123-b07; source PDF page 123
\blockanchor{p123-b07}
(1) 计算 $A = \iint_{D} \left| xy - \frac{1}{4} \right| \, \mathrm{d}x \, \mathrm{d}y$.\par
% p123-b08; source PDF page 123
\blockanchor{p123-b08}
(2) 设 $z = f(x, y)$ 在 $D$ 上连续, 且满足 $\iint_{D} f(x, y) \, \mathrm{d}x \, \mathrm{d}y = 0$, $\iint_{D} xyf(x, y) \, \mathrm{d}x \, \mathrm{d}y = 1$, 证明: 存在 $(x_0, y_0) \in D$, 使得 $Af(x_0, y_0) \geq 1$.\par
% p123-b09; source PDF page 123
\par\Needspace{6\baselineskip}
\blockanchor{p123-b09}
\question{7}(15 分) 已知二元函数\par
\nopagebreak[4]
% p123-b10; source PDF page 123
\blockanchor{p123-b10}
\displayformula{f (x, y) = \left\{ \begin{array}{l l} \frac {x ^ {3}}{x ^ {2} + y ^ {2}}, & (x, y) \neq (0, 0); \\ 0, & (x, y) = (0, 0). \end{array} \right.}
% p123-b11; source PDF page 123
\blockanchor{p123-b11}
(1) 求 $f_{x}(x, y), f_{y}(x, y)$.\par
% p123-b12; source PDF page 123
\blockanchor{p123-b12}
(2) 设方向 l 与 x 轴正方向的夹角为  $\alpha$ ，求  $\left.\frac{\partial f(x,y)}{\partial l}\right|_{(0,0)}$ .\par
% p123-b13; source PDF page 123
\blockanchor{p123-b13}
(3) 证明 $f(x, y)$ 在点 $(0, 0)$ 处连续但不可微.\par
% p123-b14; source PDF page 123
\blockanchor{p123-b14}
\question{8}(15 分) 设 $p > 1, q \in \mathbb{R}$, 证明: 函数项级数 $\sum_{n=1}^{\infty} \frac{1}{n^p + n^q x^2}$ 在 $(- \infty, +\infty)$ 上一致收敛, 且当 $q < 3p - 2$ 时可逐项微分.\par
% p123-b15; source PDF page 123
\blockanchor{p123-b15}
\question{9}(20 分) 证明下列各题.\par
% p123-b16; source PDF page 123
\blockanchor{p123-b16}
(1) (5 分) 设 $a$ 为大于 0 的常数, 当 $x > 0$ 时, $f(x) = \frac{1}{\sqrt{x}} - \frac{1}{\sqrt{x + a}}$ 单调递减.\par
% p123-b17; source PDF page 123
\blockanchor{p123-b17}
(2) (7 分) $\int_{0}^{\sqrt{2\pi}} \sin x^{2} \, \mathrm{d}x > \frac{2 - \sqrt{2}}{2\sqrt{\pi}}$.\par
% p123-b18; source PDF page 123
\blockanchor{p123-b18}
(3) (8 分) $\int_{0}^{+\infty}\sin x^{2}\mathrm{d}x$ 条件收敛.\par
% p123-b19; source PDF page 123
\blockanchor{p123-b19}
\question{10}(20 分) 设 $f(x)$ 在 $[a, b]$ 上不恒为常数, 在 $[a, b]$ 上二阶可导, 且满足 $f(a) = f(b)$, $\lim_{x \to a^{+}} \frac{f(x) - f(a)}{(x - a)^2} = 2$. 证明:\par
% p123-b20; source PDF page 123
\blockanchor{p123-b20}
(1) $f_{+}^{\prime}(a)=0,\ f_{+}^{\prime\prime}(a)=4.$\par
% p123-b21; source PDF page 123
\blockanchor{p123-b21}
(2) 存在 $\xi \in (a, b)$, 使得 $f''(\xi) = 0$.\par
% p123-b22; source PDF page 123
\blockanchor{p123-b22}
(3) 存在 $\eta \in (a, b)$, 使得 $f'(\eta) = \frac{f(\eta) - f(a)}{\eta - a}$.\par
% p123-b23; source PDF page 123
\blockanchor{p123-b23}
\question{11}(15 分) 设连续函数列 $\{f_n(x)\}$ 在 $[a, b]$ 上一致收敛于 $f(x)$, 而 $g(u)$ 在 $(- \infty, +\infty)$ 上连续, 证明: $\{g(f_n(x))\}$ 在 $[a, b]$ 上一致收敛于 $g(f(x))$.\par
